\input{preamble}

\section*{Time-independent Schrodinger equation}

Let $H(x)$ be a time-independent Hamiltonian.
Let $\psi(x)$ be a function such that $H(x)\psi(x)$ has eigenvalue $E$.
\begin{equation*}
H(x)\psi(x)=E\psi(x)
\end{equation*}

Recall time-evolution operator $U(t,t_0)$.
\begin{equation*}
U(t,t_0)=\exp\left(-\frac{i}{\hbar}(t-t_0)H(x)\right)
\end{equation*}

Let $\Psi(x,t)$ be the function
\begin{equation*}
\Psi(x,t)=U(t,0)\psi(x)=\exp\left(-\frac{i}{\hbar}tH(x)\right)\psi(x)
\end{equation*}

By hypothesis
\begin{equation*}
\Psi(x,t)=\exp\left(-\frac{i}{\hbar}Et\right)\psi(x)
\end{equation*}

By linearity of $H(x)$
\begin{equation*}
H(x)\Psi(x,t)=\exp\left(-\frac{i}{\hbar}Et\right)H(x)\psi(x)
=\exp\left(-\frac{i}{\hbar}Et\right)E\psi(x)
\tag{1}
\end{equation*}

Evaluate the Schrodinger equation time derivative.
\begin{equation*}
i\hbar\frac{\partial}{\partial t}\Psi(x,t)
=E\exp\left(-\frac{i}{\hbar}Et\right)\psi(x)
\end{equation*}

Hence by (1)
\begin{equation*}
i\hbar\frac{\partial}{\partial t}\Psi(x,t)=H(x)\Psi(x,t)
\end{equation*}

\end{document}

Let $\Psi(x,t)$ be the wave function
\begin{equation*}
\Psi(x,t)=\psi(x)\exp\left(-\frac{i}{\hbar}Et\right)
\end{equation*}

For the time derivative we have
\begin{equation*}
i\hbar\frac{d}{dt}\Psi(x,t)=E\Psi(x,t)
\tag{1}
\end{equation*}

Recall the Schrodinger equation
\begin{equation*}
i\hbar\frac{d}{dt}\Psi(x,t)=H\Psi(x,t)
\tag{2}
\end{equation*}

By equivalence of (1) and (2)
\begin{equation*}
E\Psi(x,t)=H\Psi(x,t)
\end{equation*}

Let $H$ be independent of $t$.
Then by linearity of $H$ the time exponential of $\Psi$ can be separated from the scope of $H$.
\begin{equation*}
E\Psi(x,t)=\exp\left(-\frac{i}{\hbar}Et\right)H\psi(x)
\end{equation*}

Cancel exponentials to obtain the time-independent Schrodinger equation
\begin{equation*}
E\psi(x)=H\psi(x)
\end{equation*}

\href{https://georgeweigt.github.io/examples/time-independent-schrodinger-equation-demo.html}{Eigenmath script}

\end{document}
